\documentclass[11pt, a4paper]{exam}
\usepackage[a4paper,top=2.3cm,bottom=2.8cm,left=2cm,right=2cm,headheight=15pt,headsep=0.6cm,footskip=1.2cm]{geometry}
\usepackage[T1]{fontenc}
\usepackage[utf8]{inputenc}
\usepackage{lmodern}
\usepackage{textcomp}
\usepackage[french]{babel}
\usepackage{xcolor}
\usepackage{listings}
\usepackage{graphicx}
\usepackage{array}
\usepackage{enumitem}
\usepackage{amsmath}
\usepackage{amssymb}
\usepackage{needspace}
\usepackage{tikz}
\usetikzlibrary{positioning,shapes.multipart,arrows.meta}
\usepackage[most]{tcolorbox}

\definecolor{codegreen}{rgb}{0,0.6,0}
\definecolor{codegray}{rgb}{0.5,0.5,0.5}
\definecolor{codepurple}{rgb}{0.58,0,0.82}
\definecolor{backcolour}{rgb}{0.95,0.95,0.92}

\lstdefinestyle{mystyle}{
    backgroundcolor=\color{backcolour},
    commentstyle=\color{codegreen},
    keywordstyle=\color{magenta},
    numberstyle=\tiny\color{codegray},
    stringstyle=\color{codepurple},
    basicstyle=\ttfamily\footnotesize,
    breakatwhitespace=false,
    breaklines=true,
    postbreak=\mbox{\textcolor{codegray}{$\hookrightarrow$}\space},
    captionpos=b,
    keepspaces=true,
    columns=fullflexible,
    numbers=left,
    numbersep=5pt,
    showspaces=false,
    showstringspaces=false,
    showtabs=false,
    upquote=true,
    tabsize=4,
    xleftmargin=1.5em,
    framexleftmargin=0.3em,
    aboveskip=0.6em,
    belowskip=0.6em,
    extendedchars=true,
    literate={é}{{\'e}}1 {è}{{\`e}}1 {ê}{{\^e}}1 {ë}{{\"e}}1 {à}{{\`a}}1 {â}{{\^a}}1
             {ç}{{\c c}}1 {ù}{{\`u}}1 {û}{{\^u}}1 {ô}{{\^o}}1 {î}{{\^i}}1 {ï}{{\"i}}1
             {É}{{\'E}}1 {È}{{\`E}}1 {À}{{\`A}}1 {Ç}{{\c C}}1 {œ}{{\oe}}1
             {«}{{\guillemotleft}}1 {»}{{\guillemotright}}1 {→}{{$\rightarrow$}}1
}
\lstset{style=mystyle, language=Python}
% Console : sortie d'un programme, une ligne par ligne affichée
\lstdefinestyle{console}{style=mystyle, language={}, numbers=none, backgroundcolor=\color{white}, frame=single, rulecolor=\color{codegray}, xleftmargin=0.3em}

% Code en ligne
\newcommand{\code}[1]{\texttt{#1}}

% Titre d'une question : numéro et intitulé, barème à droite
\newcommand{\Qt}[2]{\par\needspace{10\baselineskip}\subsection*{#1\hfill{\normalfont\small #2}}}
% Rappel de la consigne en une ligne
\newcommand{\rappel}[1]{\noindent\textit{Ce qui est demandé :} #1\par\smallskip}
% Barème détaillé
\newcommand{\bareme}[1]{\par\noindent\textbf{Barème.}\ #1\par\smallskip}

% La réponse attendue, mise en évidence
\newtcolorbox{reponse}[1][Réponse]{enhanced, breakable, colback=black!5, colframe=black, boxrule=0.9pt, arc=1mm,
  title={#1}, fonttitle=\bfseries\small, coltitle=black, colbacktitle=black!15, left=2mm, right=2mm, top=1mm, bottom=1mm,
  before skip=0.5em, after skip=0.7em}
% Erreur fréquente
\newtcolorbox{erreur}{enhanced, colback=white, colframe=black, boxrule=0.5pt, leftrule=3pt, arc=0mm,
  title={Erreur fréquente}, fonttitle=\bfseries\small, coltitle=black, colbacktitle=white, left=2mm, right=2mm, top=1mm, bottom=1mm,
  before skip=0.5em, after skip=0.7em}
% Encadré « À retenir »
\newtcolorbox{retenir}{enhanced, breakable, colback=white, colframe=black, boxrule=0.5pt, arc=1mm, boxsep=1mm,
  borderline={0.4pt}{1.2pt}{black, dashed},
  title={À retenir}, fonttitle=\bfseries\small, coltitle=black, colbacktitle=white, left=2mm, right=2mm, top=1mm, bottom=1mm,
  before skip=0.5em, after skip=0.7em}
% Remarque sur l'énoncé (erreur ou ambiguïté du sujet)
\newtcolorbox{remarque}{enhanced, breakable, colback=white, colframe=black, boxrule=0.5pt, arc=0mm,
  borderline={0.4pt}{1.2pt}{black, dotted},
  title={Remarque sur l'énoncé}, fonttitle=\bfseries\small, coltitle=black, colbacktitle=black!8, left=2mm, right=2mm, top=1mm, bottom=1mm,
  before skip=0.5em, after skip=0.7em}

\newcommand{\methode}{\noindent\textbf{Méthode.}\ }
\newcommand{\explication}{\noindent\textbf{Explication.}\ }

\setlist{itemsep=0.15em, topsep=0.3em}
\renewcommand{\arraystretch}{1.25}
\setlength{\parindent}{0pt}
\setlength{\parskip}{0.35em}
\setlength{\emergencystretch}{3em}

% Pile dessinée verticalement (sommet en haut) et file horizontale (tête à gauche)
\newcommand{\pilev}[2]{\begin{tikzpicture}[baseline=(titre.base)]
  \node[font=\scriptsize, align=center, text width=22mm] (titre) at (0,0) {#1};
  \foreach \v [count=\k] in {#2} {\node[draw, minimum width=9mm, minimum height=5mm, inner sep=0pt, font=\small] at (0,-0.25-\k*0.5) {\v};}
\end{tikzpicture}}
\newcommand{\fileh}[2]{\begin{tikzpicture}[baseline=(titre.base)]
  \node[font=\scriptsize, anchor=west] (titre) at (-0.45,0.55) {#1};
  \foreach \v [count=\k] in {#2} {\node[draw, minimum width=9mm, minimum height=5mm, inner sep=0pt, font=\small] at (\k*0.9-0.9,0) {\v};}
\end{tikzpicture}}


\title{Corrigé — CC Term NSI - POO, Pile, File et rappels de 1re}
\author{Babash - NSI Tani Malandi}
\date{Novembre 2026 (2 h, sur 20 points)}

\begin{document}
\pagestyle{headandfoot}
\firstpageheader{}{}{}
\runningheader{\small Corrigé}{}{\small POO, pile, file et rappels de 1re -- novembre 2026}
\runningheadrule
\firstpagefooter{}{\footnotesize Babash - NSI Tani Malandi \textperiodcentered{} CC BY-SA 4.0 \textperiodcentered{} page \thepage}{}
\runningfooter{}{\footnotesize Babash - NSI Tani Malandi \textperiodcentered{} CC BY-SA 4.0 \textperiodcentered{} page \thepage}{}

\maketitle
\thispagestyle{headandfoot}
\begin{center}\footnotesize\itshape Exercices 1 à 4 écrits pour ce contrôle. Exercice 5 : exercice 2 du sujet 24-NSIJ2JA1 (baccalauréat 2024, Asie, jour 2), adapté : questions 1, 2, 4, 5, 6 et 8 ; la question 3 remplacée par une phrase de l'énoncé, la question 7 retirée, l'en-tête des fonctions à écrire donné. Corrigé rédigé pour ce contrôle, comparé au corrigé publié par l'AEIF.\end{center}

\begin{tcolorbox}[colback=white, colframe=black, boxrule=0.5pt, arc=1mm, left=2mm, right=2mm]
\small\textbf{Comment lire ce corrigé.} Pour chaque question : la réponse attendue (encadrée), une explication si elle aide, le barème et l'erreur la plus fréquente. Essaie d'abord seul·e, puis compare ta démarche, pas seulement ton résultat. Tous les codes et toutes les sorties de ce corrigé ont été obtenus en exécutant Python 3.13 (script \code{verifier.py}).

\textbf{Pour corriger.}
\begin{itemize}[leftmargin=1.2em, itemsep=1pt, topsep=2pt]
\item Une question, une réponse attendue. Quand la réponse a plusieurs morceaux (des trous, une fonction), ses points se partagent entre eux comme le dit le barème : on cherche les points que la copie montre.
\item \textbf{Toute fonction qui fait ce qui est demandé a tous les points}, même écrite autrement ou maladroitement (une boucle de plus, un \code{if} au lieu d'un \code{return} direct).
\item \textbf{Une erreur de syntaxe mineure} dans un raisonnement juste (deux-points, parenthèse ou guillemet oublié, \code{"} au lieu de \code{'}, une faute dans un nom) : rien, ou 0,25 au plus par question.
\item \textbf{Une erreur ne se paie qu'une fois} : une réponse juste à partir d'une réponse précédente fausse a tous ses points ; un même savoir n'est noté que par une question.
\end{itemize}
\end{tcolorbox}

\begin{center}\small
\begin{tabular}{|l|l|c|c|c|}\hline
Exercice & Notion & Questions & Points & Durée conseillée \\\hline
1 & Questions de cours : POO, pile et file & 1 à 4 & 4 & 12 min \\\hline
2 & POO : lire une classe, instancier & 5 à 7 & 2,5 & 15 min \\\hline
3 & Pile et file & 8 et 9 & 2,5 & 15 min \\\hline
4 & Rappels de 1re : dictionnaire & 10 & 1 & 8 min \\\hline
5 & Bac Asie 2024 j2 ex 2 : parenthèses et pile & 11 à 16 & 10 & 55 min \\\hline
 & Total (et 15 min de relecture) & 23 réponses & 20 & 2 h \\\hline
\end{tabular}
\end{center}

% =====================================================================
\needspace{16\baselineskip}
\section*{Exercice 1 -- Questions de cours (4 points)}

\Qt{Question 1 -- LIFO}{0,5 point}
\begin{reponse}
\textit{Last In, First Out} : dernier entré, premier sorti.
\end{reponse}
\bareme{0,5 en anglais ou en français.}

\Qt{Question 2 -- Le paramètre \code{self}}{0,5 point}
\begin{reponse}
\textbf{b)} l'objet sur lequel la méthode est appelée (dans \code{b1.places\_libres()}, \code{self} est \code{b1}).
\end{reponse}
\bareme{0,5 ; 0 si deux lettres sont données.}
\begin{erreur}
a) « la classe dans laquelle la méthode est définie » : \code{self} est un objet, une instance de la classe, pas la classe elle-même.
\end{erreur}

\Qt{Question 3 -- Vrai ou faux}{4 × 0,5 point}
\begin{reponse}
a) \textbf{V} : Python appelle \code{\_\_init\_\_} à chaque création d'objet, par exemple \code{Barge("Corail", 400)}.\\
b) \textbf{F} : \code{sommet(p)} lit l'élément du sommet sans le retirer ; c'est \code{depiler(p)} qui le retire.\\
c) \textbf{F} : on crée autant d'objets qu'on veut à partir d'une classe (\code{b1}, \code{b}…).\\
d) \textbf{V} : par exemple \code{embarquer} modifie \code{self.passagers}.
\end{reponse}
\bareme{0,5 par affirmation.}

\Qt{Question 4 -- Associer situations et structures}{1 point}
\begin{reponse}
a : \textbf{file} (premier envoyé, premier imprimé) ; b : \textbf{pile} (dernière page visitée, première revue) ; c : \textbf{pile} (dernière action faite, première annulée) ; d : \textbf{file} (premier arrivé, premier servi).
\end{reponse}
\bareme{0,25 par situation.}

% =====================================================================
\needspace{16\baselineskip}
\section*{Exercice 2 -- POO : lire une classe et créer des objets (2,5 points)}

\begin{retenir}
Une \textbf{classe} est un modèle ; un \textbf{objet} (ou \textbf{instance}) est fabriqué à partir de ce modèle. Les \textbf{attributs} sont les données de l'objet (\code{self.nom}…) ; les \textbf{méthodes} sont les fonctions définies dans la classe. On crée un objet en appelant la classe : \code{Barge("Corail", 400)} ; Python appelle alors \code{\_\_init\_\_}.
\end{retenir}

\Qt{Question 5 -- Les attributs}{0,5 point}
\begin{reponse}
\code{nom}, \code{capacite} et \code{passagers}.
\end{reponse}
\bareme{0,5 pour les trois ; 0,25 s'il en manque un.}
\begin{erreur}
Oublier \code{passagers}, qui est bien un attribut (\code{self.passagers}) même s'il n'est pas un paramètre de \code{\_\_init\_\_} ; ou citer \code{n}, qui n'est qu'un paramètre de \code{embarquer}.
\end{erreur}

\Qt{Question 6 a) -- Créer une barge}{0,5 point}
\begin{reponse}
\code{b1 = Barge("Corail", 400)}
\end{reponse}
\bareme{0,5 ; 0,25 si l'objet n'est pas stocké dans \code{b1}.}
\begin{erreur}
\code{b1 = Barge(self, "Corail", 400)} : on n'écrit jamais \code{self} à l'appel, Python le passe tout seul.
\end{erreur}

\Qt{Question 6 b) -- Faire embarquer}{0,5 point}
\begin{reponse}
\code{b1.embarquer(120)}
\end{reponse}
\bareme{0,5 ; on accepte \code{Barge.embarquer(b1, 120)}. \code{b1.passagers = b1.passagers + 120} : 0,25, car il ne vérifie pas les places libres.}
\begin{erreur}
\code{embarquer(b1, 120)} sans le point : \code{embarquer} n'existe pas en dehors de la classe (\code{NameError: name 'embarquer' is not defined}).
\end{erreur}

\Qt{Question 7 a) et b) -- Prévoir un affichage}{0,5 + 0,5 point}
\begin{reponse}
a) La ligne 4 affiche \code{50}.\qquad b) La ligne 7 affiche \code{230}.
\end{reponse}
\explication
\begin{center}\small
\begin{tabular}{|c|l|c|c|}\hline
Ligne & Instruction & \code{passagers} & places libres \\\hline
1 & \code{b = Barge("Lagon", 300)} & 0 & 300 \\\hline
2 & \code{b.embarquer(250)} & 250 & 50 \\\hline
3 & \code{b.embarquer(80)} : refusé (80 $>$ 50) & 250 & 50 \\\hline
5 & \code{b.passagers = b.passagers - 100} & 150 & 150 \\\hline
6 & \code{b.embarquer(80)} : accepté & 230 & 70 \\\hline
\end{tabular}
\end{center}
\bareme{0,5 chacune ; une ligne 7 juste à partir d'une ligne 4 fausse a ses 0,5.}
\begin{erreur}
Ajouter 80 à la ligne 3 alors que l'embarquement est refusé (on trouve alors $-30$ à la ligne 4) : quand la condition est fausse, la méthode renvoie \code{False} sans modifier \code{passagers}.
\end{erreur}

% =====================================================================
\needspace{16\baselineskip}
\section*{Exercice 3 -- Pile et file (2,5 points)}

\Qt{Question 8 a) -- La pile après la boucle}{0,5 point}
\begin{reponse}
\pilev{pile \code{p}\\(sommet en haut)}{9,1,7,4}
\end{reponse}
\bareme{0,5 ; 0 si la pile est à l'envers (4 en haut). Un dessin horizontal où le sommet est clairement désigné est accepté.}

\Qt{Question 8 b) -- Prévoir un affichage}{0,5 point}
\begin{reponse}
\begin{lstlisting}[style=console]
13 7 4 4
\end{lstlisting}
\end{reponse}
\explication \code{a} vaut 9 (le dernier empilé), \code{c} vaut 4 (le premier enfilé). On empile \code{a + c}, soit 13, puis on enfile le sommet de \code{p}, soit 13. À la fin, \code{p} contient (du sommet au fond) 13, 1, 7, 4 et \code{f} contient (de la tête à la queue) 7, 1, 9, 13.
\bareme{0,5 pour les quatre valeurs dans l'ordre ; 0,25 pour trois justes.}
\begin{erreur}
Croire que \code{sommet(p)} retire l'élément : il le lit seulement, donc \code{taille(p)} vaut toujours 4.
\end{erreur}

\Qt{Question 9 a) -- Les lignes de \code{inverser}}{1 point}
\begin{reponse}
\begin{lstlisting}[firstnumber=4]
        empiler(p, defiler(f))
\end{lstlisting}
\begin{lstlisting}[firstnumber=6]
    while not est_vide(p):
        enfiler(g, depiler(p))
    return g
\end{lstlisting}
\end{reponse}
\bareme{0,25 par ligne juste. On ne demande pas de recopier le reste du code ; une ligne sans son numéro mais reconnaissable est acceptée.}
\begin{erreur}
\code{empiler(p, tete(f))} : \code{tete} ne retire rien, la boucle \code{while} ne s'arrête jamais. Ou \code{return p} : on renvoie la pile, pas une file.
\end{erreur}

\Qt{Question 9 b) -- La file \code{f} après l'appel}{0,5 point}
\begin{reponse}
Rien : la file \code{f} est \textbf{vide}.
\end{reponse}
\explication La première boucle vide \code{f} en la défilant, et la fonction travaille sur la file passée en paramètre (elle ne la copie pas). La réponse ne dépend pas de la question 9 a) : l'énoncé dit ce que fait la fonction ; la boucle \code{while not est\_vide(f)}, donnée, suffit.
\bareme{0,5.}
\begin{erreur}
Croire que \code{f} n'a pas changé : défiler retire vraiment les éléments de la file passée en paramètre.
\end{erreur}

% =====================================================================
\needspace{16\baselineskip}
\section*{Exercice 4 -- Rappels de 1re : dictionnaire (1 point)}

\Qt{Question 10 -- Trouver l'erreur}{1 point (0,5 + 0,5)}
\begin{reponse}
Ligne 5 corrigée : \code{    for produit in panier:}\\
Explication : la boucle parcourt les produits de \code{prix} ; au deuxième tour, \code{produit} vaut \code{"bananes"}, qui n'est pas une clé du panier, donc \code{panier[produit]} provoque l'erreur \code{KeyError: 'bananes'}. Il faut parcourir le panier : chacun de ses produits a un prix.
\end{reponse}
\explication Avec la ligne corrigée, \code{cout(\{"mangues": 2, "tomates": 0.5\}, prix)} renvoie \code{7.0} (2 × 2,5 + 0,5 × 4).
\bareme{0,5 pour la ligne corrigée (toute ligne 5 qui fait renvoyer \code{7.0}, par exemple \code{for produit in panier.keys():}) ; 0,5 pour l'explication : on parcourt les produits de \code{prix}, dont \code{"bananes"}, qui n'est pas une clé de \code{panier}. Une explication juste sans correction a ses 0,5.}
\begin{erreur}
Croire que l'erreur est à la ligne 6, désignée par le message : le message montre où le programme s'arrête, la cause est le choix du dictionnaire parcouru, ligne 5.
\end{erreur}

% =====================================================================
\needspace{16\baselineskip}
\section*{Exercice 5 -- Type bac : parenthèses et pile (10 points)}
\begin{center}\footnotesize\itshape Sujet 24-NSIJ2JA1, baccalauréat 2024, Asie, jour 2, exercice 2 (6 points au bac), adapté comme le dit la tête du sujet. Le barème par question est celui de ce contrôle.\end{center}

\Qt{Question 11 (question 1 du bac) -- L'expression est-elle bien parenthésée ?}{1,5 point}
\begin{reponse}
\textbf{Non.} On lit \code{[} puis \code{(} : les ouvrantes encore ouvertes sont \code{[ (}. Le \code{)} après \code{i+1} ferme \code{(} : il reste \code{[}. Le \code{)} après \code{-3} devrait fermer la dernière ouvrante encore ouverte, qui est \code{[} : les deux ne forment pas un couple.
\end{reponse}
\explication On peut aussi remarquer que l'expression contient 3 ouvrantes et 4 fermantes ; cet argument suffit ici, mais il ne suffit pas en général (question 14).
\bareme{La question du bac demande de répondre et de justifier : 0,5 pour « non » ; 1 pour une justification qui désigne la parenthèse fautive (le \code{)} après \code{-3}) ou le décompte 3 ouvrantes, 4 fermantes.}

\Qt{Question 12 (question 2 du bac) -- \code{compte\_ouvrante}}{1,5 point}
\begin{reponse}
\begin{lstlisting}
def compte_ouvrante(txt):
    n = 0
    for c in txt:
        if c in "([{":
            n = n + 1
    return n
\end{lstlisting}
\end{reponse}
\explication \code{c in "([\{"} vaut \code{True} si le caractère \code{c} est l'un des trois caractères de la chaîne. On accepte \code{if c == "(" or c == "[" or c == "\{":}. L'en-tête (signature et docstring) n'a pas à être recopié.
\bareme{1,5 pour toute fonction juste. Sinon : 0,5 parcours de la chaîne avec un compteur ; 0,5 test des trois types ; 0,5 \code{return} après la boucle. Un seul type de parenthèse testé : 1 au plus.}
\begin{erreur}
\code{if c == "(" or "[" or "\{":} : la condition est toujours vraie (\code{"["} est une chaîne non vide), on compte tous les caractères. Il faut répéter \code{c ==} à chaque fois.
\end{erreur}

\Qt{Question 13 (question 4 du bac) -- \code{bon\_compte}}{1 point}
\begin{reponse}
\begin{lstlisting}
def bon_compte(txt):
    return compte_ouvrante(txt) == compte_fermante(txt)
\end{lstlisting}
\end{reponse}
\explication La comparaison \code{==} vaut déjà \code{True} ou \code{False} : on peut la renvoyer directement. La version avec \code{if \ldots: return True else: return False} est juste aussi (tous les points). La question ne dépend pas de la question 12 : on utilise les deux fonctions sans les réécrire (la phrase de l'énoncé donne \code{compte\_fermante}).
\bareme{1 pour toute fonction juste ; 0,5 si l'on recompte les parenthèses sans utiliser les deux fonctions, comme l'énoncé le demande.}

\Qt{Question 14 (question 5 du bac) -- Un contre-exemple}{1 point}
\begin{reponse}
Par exemple \code{')('} : une ouvrante et une fermante, donc \code{bon\_compte} renvoie \code{True}, mais la fermante arrive avant l'ouvrante. Autres exemples : \code{'(]'}, \code{'([)]'}.
\end{reponse}
\bareme{1 pour un exemple juste (vérifié par le script).}
\begin{erreur}
\code{'tab[2*(i + 4] - tab[3)'} n'est pas un bon exemple : il a 3 ouvrantes et 2 fermantes, \code{bon\_compte} renvoie \code{False}.
\end{erreur}

\Qt{Question 15 (question 6 du bac) -- Lignes 9 et 14 de la classe \code{Pile}}{1,5 point}
\begin{reponse}
\begin{lstlisting}
    def empiler(self, elt):
        self.contenu.append(elt)

    def depiler(self):
        if self.est_vide():
            return "La pile est vide."
        return self.contenu.pop()
\end{lstlisting}
\end{reponse}
\explication Le sommet de la pile est la fin de la liste \code{self.contenu} : \code{append} ajoute à la fin, \code{pop()} retire le dernier élément \textbf{et le renvoie}.
\bareme{ligne 9 : 0,75 ; ligne 14 : 0,75. On accepte \code{self.contenu.insert(0, elt)} avec \code{self.contenu.pop(0)} (sommet au début), si les deux sont cohérents.}
\begin{erreur}
\code{self.contenu.pop(0)} avec \code{append} : on retire le premier entré, c'est une file. Ou \code{return self.contenu[-1]} : on lit le sommet sans le retirer.
\end{erreur}

\Qt{Question 16 (question 8 du bac) -- \code{est\_bien\_parenthesee}}{3,5 points}
\begin{reponse}
\begin{lstlisting}
def est_bien_parenthesee(txt):
    couples = {")": "(", "]": "[", "}": "{"}
    p = Pile()
    for c in txt:
        if c in "([{":
            p.empiler(c)
        elif c in ")]}":
            if p.depiler() != couples[c]:
                return False
    return p.est_vide()
\end{lstlisting}
\end{reponse}
\explication
\begin{itemize}
\item Ligne 2 : le dictionnaire associe à chaque fermante son ouvrante. Sans dictionnaire, on accepte un test comme \code{(o == "(" and c == ")") or (o == "[" and c == "]") or (o == "\{" and c == "\}")}.
\item Ligne 8 : si la pile est vide, \code{depiler} renvoie la chaîne \code{"La pile est vide."}, qui n'est égale à aucune ouvrante : la fonction renvoie bien \code{False}, sans erreur.
\item Ligne 10 : à la fin, il ne doit rester aucune ouvrante non fermée.
\end{itemize}
Résultats obtenus en exécutant le code :
\begin{center}\small
\begin{tabular}{|l|c|}\hline
Appel & Valeur renvoyée \\\hline
\code{est\_bien\_parenthesee('tab[2*(i + 4)] - tab[3]')} & \code{True} \\\hline
\code{est\_bien\_parenthesee('tab[2*(i + 4] - tab[3)')} & \code{False} \\\hline
\code{est\_bien\_parenthesee('[2*(i+1)-3) for i in range(3, 10)]')} & \code{False} \\\hline
\code{est\_bien\_parenthesee('((')} & \code{False} \\\hline
\code{est\_bien\_parenthesee(')')} & \code{False} \\\hline
\end{tabular}
\end{center}
\bareme{3,5 pour toute fonction juste, même construite autrement. Sinon, on cherche les points que la copie montre : 0,5 création de la pile ; 0,5 parcours ; 0,5 empiler les ouvrantes ; 0,5 dépiler sur une fermante ; 1 test du couple et \code{return False} ; 0,5 \code{return p.est\_vide()} après la boucle.}
\begin{erreur}
\code{if p.depiler() != c:} : on compare \code{"("} à \code{")"}, toujours différents ; il faut comparer l'ouvrante dépilée à l'ouvrante qui correspond à \code{c}. Autre erreur : \code{return True} à la fin sans tester la pile (\code{'(('} serait déclarée bien parenthésée).
\end{erreur}

\begin{retenir}
Une pile sert à vérifier ce qui s'ouvre et se ferme dans l'ordre inverse : la dernière parenthèse ouverte est la première à fermer (LIFO). Le simple comptage ne voit pas l'ordre.
\end{retenir}
\end{document}
